%\documentclass[12pt,letterpaper]{article}



%%%
%% Required packages:
%%
\usepackage{etex}
\usepackage{amsmath}
\usepackage{amssymb}
\usepackage{amstext}
\usepackage{amsthm}
\usepackage{fancyhdr,fancybox}
\usepackage{mathrsfs} % need for \mathscr command
\usepackage{multicol}
\usepackage[normalem]{ulem}
%\usepackage{pictex,pictexplus}
% \usepackage{pictexwd}
\usepackage{rotating}
\usepackage{url}
\usepackage{xfrac}
\usepackage{booktabs}
\usepackage{color,soul}
\usepackage{comment}

\usepackage{hyperref}
\hypersetup{
    colorlinks=true,     % false: boxed links; true: colored links
    linkcolor=blue,      % color of internal links
    citecolor=blue,      % color of links to bibliography
    filecolor=blue,      % color of file links
    urlcolor=blue        % color of external links
}


\usepackage{tikz}
\usetikzlibrary{backgrounds,calc}

%%
%% Common formatting commands
%%
\setlength{\parindent}{0in}
\setlength{\textwidth}{7in}
\setlength{\evensidemargin}{-0.25in}
\setlength{\oddsidemargin}{-0.25in}
\setlength{\parskip}{.5\baselineskip}
\setlength{\topmargin}{-0.5in}
\setlength{\textheight}{9in}

%               Problem and Part
\newcounter{problemnumber}
\newcounter{partnumber}
\newcounter{subpartnumber}
\newcommand{\Problem}{%
  \refstepcounter{problemnumber}%
  \setcounter{partnumber}{0}%
  \setcounter{subpartnumber}{0}%
  \item[\makebox{\hfil\textbf{\theproblemnumber.}\hfil}]%
  }
\newcommand{\InProblem}[2][1.6]{%
  \refstepcounter{problemnumber}%
  \setcounter{partnumber}{0}%
  \setcounter{subpartnumber}{0}%
  \textbf{\theproblemnumber.}\ \ \parbox[t]{#1in}{#2}%
}
\newcommand{\InSmallProblem}[1]{\InProblem[1.05]{#1}}
\newcommand{\NoProblem}[1]{%
  \mbox{\hfil\textbf{#1}\hfil}%
  }
\newcommand{\UnProblem}{%
  \refstepcounter{problemnumber}%
  \setcounter{partnumber}{0}%
  \setcounter{subpartnumber}{0}%
  \mbox{\hfil\textbf{\theproblemnumber.}\hfil}%
  }
\newcommand{\InFlexProblem}[1]{%
  \refstepcounter{problemnumber}%
  \setcounter{partnumber}{0}%
  \setcounter{subpartnumber}{0}%
  \textbf{\theproblemnumber.}\ \ #1%
}
%%  Part:
\newcommand{\Part}{%
  \renewcommand{\thepartnumber}{(\alph{partnumber})}%
  \refstepcounter{partnumber}
  \setcounter{subpartnumber}{0}%
  \item[(\alph{partnumber})]%
}
\newcommand{\InPart}[2][1.75]{%
  \renewcommand{\thepartnumber}{(\alph{partnumber})}%
  \refstepcounter{partnumber}%
  \setcounter{subpartnumber}{0}%
  (\alph{partnumber})\ \ \parbox[t]{#1in}{#2}%
}
\newcommand{\UnPart}{%
  \refstepcounter{partnumber}%
  \renewcommand{\thepartnumber}{(\alph{partnumber})}%
  \setcounter{subpartnumber}{0}%
  (\alph{partnumber})%
}
\newcommand{\InLargePart}[1]{\InPart[3.05]{#1}}
\newcommand{\InFullPart}[1]{\InPart[6.2]{#1}}
\newcommand{\InSmallPart}[1]{\InPart[1.05]{#1}}
%% SubPart:
\newcommand{\SubPart}{%
  \refstepcounter{subpartnumber}%
  \item[(\roman{subpartnumber})]%
  }
\newcommand{\InSubPart}[2][2.25]{%
  \stepcounter{subpartnumber}%
  (\roman{subpartnumber})\ \ \parbox[t]{#1in}{#2}%
}
\newcommand{\InSmallSubPart}[1]{\InSubPart[1.5]{#1}}
\newcommand{\InTinySubPart}[1]{%
  \stepcounter{subpartnumber}%
  (\roman{subpartnumber})\ \ \makebox{#1}%
}
\newcommand{\InMySubPart}[2][2.25]{%
  \stepcounter{subpartnumber}%
  \parbox[t]{#1in}{\makebox[0.3in][l]{(\roman{subpartnumber})}\ \ {#2}}%
}
\newcommand{\TrueFalse}{\textbf{T}\ \ \textbf{F}\ \ }
\newcommand{\TFPart}[1]{% 
  \stepcounter{partnumber}%
  \setcounter{subpartnumber}{0}%
  \item[(\alph{partnumber})] \TrueFalse\parbox[t]{5.5in}{#1}}

\newcommand{\Points}[1]{(#1 points) \ }
\newcommand{\Point}[1]{(#1 point) \ }


%% For Answers:
\newcounter{answernumber}
\newcommand{\Answer}{%
  \refstepcounter{answernumber}%
  \setcounter{partnumber}{0}%
  \setcounter{subpartnumber}{0}%
  \item[\makebox{\hfil\textbf{\theanswernumber.}\hfil}]%
  }


% Setting up the headers for the first page
%\chead{} % will default to what is typed in the file.
\fancyhead[L]{\textsc{Math 22a}}
\fancyhead[R]{\textsc{Fall 2024}}
\fancyfoot[R]{
	\vspace*{\fill}\footnotesize\textcopyright{} \emph{Harvard Math 22a}	
}
\renewcommand{\headrulewidth}{0pt}

%\fancyhead[C]{}
%	\chead{}	
%	\lhead{}
%	\rhead{}
%	\cfoot{}


% Putting copyright on the footer of each page
%\fancypagestyle{secondpage}{
%	\fancyhead[C]{TESTing again} % change this to be blank
%		%\chead{}	
%	\fancyhead[L]{bears} % change this to be blank
%	\fancyhead[R]{dogs} % change this to be blank
%	\fancyfoot[R]{ %keeps the footer the same
%		\vspace*{\fill}\footnotesize\textcopyright{} \emph{Harvard Math 22a}	
%	}
%}
%\renewcommand{\headrulewidth}{0pt}
%\pagestyle{fancy}
\pagestyle{empty}


%\lhead{}
%\rhead{}
%\rfoot{\vspace*{\fill}\footnotesize\textcopyright{} \emph{Harvard Math 22a}}
%\renewcommand{\headrulewidth}{0pt}
%\pagestyle{fancy}




%%
%% Common shortcut commands:
%%
\newcommand{\ds}{\displaystyle}
\newcommand{\sgn}{\operatorname{sgn}}
\renewcommand{\Pr}{\operatorname{P}}
\newcommand{\CPr}[3][{}]{\Pr_{\text{#1}}\left(#2\,|\,#3\right)}

\newcommand{\C}{\mathbb{C}}
\newcommand{\R}{\mathbb{R}}
\newcommand{\Z}{\mathbb{Z}}
\newcommand{\N}{\mathbb{N}}
\newcommand{\Q}{\mathbb{Q}}
\newcommand{\D}{\mathbb{D}}

\newcommand{\rref}{\operatorname{rref}}
\newcommand{\im}{\operatorname{im}}
\newcommand{\rank}{\operatorname{rank}}
\newcommand{\diag}{\operatorname{diag}}
\newcommand{\Span}{\operatorname{span}}
%\renewcommand{\vec}[1]{\mathbf{#1}}
\newcommand{\charpoly}{\mathscr{P}}
\newcommand{\nicefrac}[2]{\sfrac{#1}{#2}} % using xfrac package
\newcommand{\topology}{\mathcal{T}}
\newcommand{\Inn}{\operatorname{Inn}}

\newcommand{\blankbox}[2]{\fbox{\rule{#1}{0in}\rule{0in}{#2}}}

\newenvironment{myindentpar}[1]%
 {\begin{list}{}%
         {\setlength{\leftmargin}{#1}
          \setlength{\rightmargin}{\leftmargin}}%
         \item[]%
 }
 {\end{list}}
\newtheorem{theorem}{Theorem}
\newtheorem*{NStheorem}{Theorem}
\newtheorem{cor}[theorem]{Corollary}
\newtheorem*{claim}{Claim}
\newtheorem{defn}{Definition}

\newenvironment{amatrix}[1]{%
  \left(\begin{array}{@{}*{#1}{c}|c@{}}
}{%
  \end{array}\right)
}

%--------grstep
% For denoting a Gauss' reduction step.
% Use as: \grstep{\rho_1+\rho_3} or \grstep[2\rho_5 \\ 3\rho_6]{\rho_1+\rho_3}
\newcommand{\grstep}[2][\relax]{%
   \ensuremath{\mathrel{
       {\mathop{\longrightarrow}\limits^{#2\mathstrut}_{
                                     \begin{subarray}{l} #1 \end{subarray}}}}}}
\newcommand{\swap}{\leftrightarrow}




\section{{Row Reduction and Row Echelon Forms, $\vec\S$1.2,  1.3}}% 
\chead{\large\textbf{Row Reduction and Row Echelon Forms, $\vec\S$1.2,  1.3}}% 

%\begin{document}
\thispagestyle{fancy}

Recall from last week:

\begin{itemize}

\item A {\bf linear equation} in the variable $x_1, \dots, x_n$ is an equation that can be put in the following form
\begin{equation*}
a_1 x_1 +a_2 x_2 +\dots + a_n x_n =b
\end{equation*}
where $b$ and the coefficients $a_1,\dots, a_n$ are real or complex numbers. A {\bf system of linear equations} is a collection of one or more linear equations involving the same variables. A {\bf solution} to the linear system is a list of $(s_1, \dots, s_n)$ of numbers that are solutions to each of the linear equations in the system. 

\item A linear system is {\bf consistent} if there is a solution and {\bf inconsistent} if there is no solution.

\item A linear system has either no solution, exactly one solution, or infinitely many solutions.

\item The following elementary row operations preserve solutions.
\begin{enumerate}
\item Replace one row by the sum of itself and a multiple of another row.
\item Interchange two rows.
\item Multiply all entries in a row by a non-zero constant.
\end{enumerate}

\item Two matrices are called {\bf row equivalent} if there is a sequence of elementary row operations changing one matrix to another.

\item If the augmented matrices of two linear systems are row equivalent, then they both have the same solution set. 

\item Proof Techniques
\begin{enumerate}
\item Quantifiers and counterexamples.
\item Direct Proofs.
\item Proof by cases.
\item Proof by contrapositive.
\end{enumerate}


\end{itemize}

\newpage

\begin{defn}
A matrix is in {\bf row echelon form} if:
\begin{enumerate}
\item All rows consisting of only zeros are at the bottom.
\item The leading coefficient (also called the pivot) of a non-zero row is always strictly to the right of the leading coefficient of the row above it.
\end{enumerate}
\end{defn}

\begin{enumerate}
  \Problem 
  \Part Give an example of a matrix in echelon form.
  
  \vfill
  
  \Part Give an example of a matrix not in echelon form.

\vfill

\begin{defn}
A matrix in row echelon form is in {\bf reduced row echelon form} if 
\begin{enumerate}
\item The leading entries are all 1.
\item Each leading 1 is the only non-zero entry in its column.
\end{enumerate}
\end{defn}

\Part Give an example in reduced row echelon form.

\vfill

\newpage

{\bf Fact:} Each matrix is row equivalent to one and only one reduced echelon matrix.

\begin{defn}
A {\bf pivot position} in a matrix $A$ is a location in $A$ that corresponds to a leading 1 in the reduced row echelon form. A {\bf pivot column} is a column that has a pivot position. 
\end{defn}

{\bf Row Reduction Algorithm (Gaussian Elimination).}
\begin{enumerate}
\item The leftmost non-zero column is a pivot column.
\item Select a non-zero entry in the pivot column as pivot. Interchange rows to move this entry to the top.
\item Use elementary row operations to create zeros below the pivot.
\item Apply (a)-(c) to the submatrix obtained by removing pivot row and column. %Repeat until all non-zero rows are modified.
Repeat until the matrix is in row echelon form.
\item Beginning with rightmost pivot and working upwards and to the left, create 0's above each pivot. Scale pivot positions entries to be 1.
\end{enumerate}

\Problem Row-reduce the following matrix:
$$\begin{bmatrix}
1 & 2 & 4 & 8\\
2 & 4 & 6 & 8\\
3 & 6 & 9 & 12\\
\end{bmatrix}$$

\vfill 

\Problem Use row-reduction to solve the linear system
$$\begin{cases}
x_1+3x_2+4x_3 =7&\\
3x_1 +9x_2 +7x_3=6&
\end{cases}$$

\vfill

\newpage

\begin{defn}
The variables corresponding to the pivot columns in an augmented matrix are called {\bf basic} and the others are called {\bf free}.
\end{defn}




{\bf Facts:} A linear system is consistent if and only if the rightmost column in the augmented matrix is not a pivot column. If the system is consistent, then it has a unique solution when there are no free variables and infinite solutions when there is one or more free variables.

\Problem Find the general solution of the augmented matrix :
$$\begin{amatrix}{3}
3 & -2 & 4 & 0\\
9& -6& 12 & 0\\
6& -4 & 8 &0\\
\end{amatrix}$$

\newpage

\begin{defn}
A matrix with 1 column is called a {\bf column vector}. A matrix with one row is called a {\bf row vector}.
\end{defn}

We can add vectors of the same size by adding the corresponding components; namely,
if $\vec{u}=\begin{bmatrix} u_1\\ \vdots\\ u_n \end{bmatrix}$ and  $\vec{v}=\begin{bmatrix} v_1\\ \vdots\\ v_n \end{bmatrix}$, then $\vec{u}+\vec{v} =\begin{bmatrix} u_1+v_1\\ \vdots\\ u_n+v_n \end{bmatrix}$. Similarly, we can scale a vector by a real number $c$, by multiplying all the components by $c$; namely, $c\vec{u}=\begin{bmatrix} cu_1\\ \vdots\\ cu_n \end{bmatrix}$.

\begin{defn} Given $\vec{v}_1,\dots, \vec{v}_p \in \mathbb{R}^n$ and scalars $c_1, \dots c_p$, the vector $$\vec{y} = c_1 \vec{v}_1 +\dots c_p \vec{v}_p$$ is called the {\bf linear combination} of $\vec{v}_1,\dots, \vec{v}_p$ with weights $c_1,\dots,c_p$.
\end{defn}

\Problem Suppose $\vec{u} =\begin{bmatrix} 1\\1\\ \end{bmatrix}$;  $\vec{v} =\begin{bmatrix} 2\\2\\ \end{bmatrix}$; and  $\vec{w} =\begin{bmatrix} 5\\2\\ \end{bmatrix}$. Can $\vec{w}$ be expressed as a linear combination of $\vec{u}$ and $\vec{v}$?


\newpage

\Problem Suppose $\vec{a} =\begin{bmatrix} 1\\-1\\ 0\\ \end{bmatrix}$; $\vec{b} =\begin{bmatrix} 0\\1\\ 2\\ \end{bmatrix}$; $\vec{c} =\begin{bmatrix} 5\\-6\\8\\ \end{bmatrix}$; and  $\vec{d} =\begin{bmatrix} 2\\-1\\6\\ \end{bmatrix}$. Can $\vec{d}$ be expressed as a linear combination of $\vec{a}$, $\vec{b}$, and $\vec{c}$?


\newpage

A vector equation $$x_1 \vec{a}_1 + x_2 \vec{a}_2 +\dots + x_n \vec{a}_n =\vec{b}$$ has the same solutions as the augment matrix $$[\vec{a}_1 \; \vec{a}_2 \; \dots \; \vec{a}_n \Big| \vec{b}].$$ Hence $\vec{b}$ is a linear combination of $\vec{a}_1, \dots, \vec{a}_n$ if and only if there is a solution to this matrix.

\begin{defn}
If $\vec{v}_1, \dots, \vec{v}_p \in \mathbb{R}^n$, then the set of all linear combinations of $\vec{v}_1, \dots, \vec{v}_p$ is called the {\bf span} of $\vec{v}_1, \dots, \vec{v}_p$. We write this set as $\text{Span} \{\vec{v}_1, \dots, \vec{v}_p\}$.
\end{defn}

\Problem 
\Part Suppose $\vec{u} =\begin{bmatrix} 1\\2\\ \end{bmatrix}$;  $\vec{v} =\begin{bmatrix} 3\\4\\ \end{bmatrix}$; and  $\vec{b} =\begin{bmatrix} -1\\4\\ \end{bmatrix}$. Is $\vec{b}$ in $\text{Span} \{ \vec{u}, \vec{v} \}$?

\vfill

\Part Does there exist $\vec{b}$ not in $\text{Span} \{ \vec{u}, \vec{v} \}$?

\vfill

\end{enumerate}

% solutions start below
%\end{document}
\begin{comment}

\newpage
\chead{\Large\textbf{Row Echelon Form and Vector Equations -- Answers / Solutions}}%
\lhead{}
\rhead{}
\thispagestyle{fancy}

\begin{enumerate}
  \Answer (There are many examples for each of these.)
  \begin{enumerate}
  \item
    \[
      \left[
        \begin{array}{cccc}
          0&1&10&\pi\\
          0&0&0&\ln(2)\\
          0&0&0&0
        \end{array}
      \right]
    \]
  \item
    \[
      \left[
        \begin{array}{cccc}
          0&0&0&1\\
          0&0&1&0\\
          0&1&0&0
        \end{array}
      \right]
    \]
  \item
    \[
      \left[
        \begin{array}{cccc}
          0&1&10&0\\
          0&0&0&1\\
          0&0&0&0
        \end{array}
      \right]
    \]
  \end{enumerate}

  \Answer

  Follow the algorithm:

  The left column is nonzero, so it is a pivot column. The top entry $1$ is not zero, so we use this for our pivot. Then we create zeroes below the pivot:

  \[
    \left[
      \begin{array}{cccc}
        1&2&4&8\\
        2&4&6&8\\
        3&6&9&12
      \end{array}
    \right]\xrightarrow{R_2+(-2)R_1\rightarrow R_2}
    \left[
      \begin{array}{cccc}
        1&2&4&8\\
        0&0&-2&-8\\
        3&6&9&12
      \end{array}
    \right]\xrightarrow{R_3+(-3)R_1\rightarrow R_3}
    \left[
      \begin{array}{cccc}
        1&2&4&8\\
        0&0&-2&-8\\
        0&0&-3&-12
      \end{array}
    \right]
  \]

  Now we focus on the submatrix
  \[
    \left[
      \begin{array}{ccc}
        0&-2&-8\\
        0&-3&-12
      \end{array}
    \right]
  \]
  The leftmost nonzero column is the second column, so it is a pivot column. The top entry $-2$ is not zero, so we use this for our pivot. Then we create zeroes below the pivot:
  \[
    \left[
      \begin{array}{ccc}
        0&-2&-8\\
        0&-3&-12
      \end{array}
    \right]\xrightarrow{R_2+(-\frac{3}{2})R_1\rightarrow R_2}
    \left[
      \begin{array}{ccc}
        0&-2&-8\\
        0&0&0
      \end{array}
    \right]
  \]
  Putting this back into the original matrix:
  \[
    \left[
      \begin{array}{cccc}
        1&2&4&8\\
        0&0&-2&-8\\
        0&0&0&0
      \end{array}
    \right]
  \]
  Now all nonzero rows have been modified, so we need to zero entries above pivots starting from the right, and then scale the pivots to be $1$.
  \[
    \left[
      \begin{array}{cccc}
        1&2&4&8\\
        0&0&-2&-8\\
        0&0&0&0
      \end{array}
    \right]\xrightarrow{R_1+2R_2\rightarrow R_2}
    \left[
      \begin{array}{cccc}
        1&2&0&-8\\
        0&0&-2&-8\\
        0&0&0&0
      \end{array}
    \right]\xrightarrow{(-\frac{1}{2})R_2\rightarrow R_2}
    \left[
      \begin{array}{cccc}
        1&2&0&-8\\
        0&0&1&4\\
        0&0&0&0
      \end{array}
    \right]
  \]

  \Answer Write this as an augmented matrix:
  \[
    \left[
      \begin{array}{ccc|c}
        1&3&4&7\\
        3&9&7&6
      \end{array}
    \right]
  \]
  and follow the algorithm again. The left column is nonzero, so it is a pivot column. The top entry $1$ is not zero, so we use this for our pivot. Then we create zeroes below the pivot. 
  \[
    \left[
      \begin{array}{ccc|c}
        1&3&4&7\\
        3&9&7&6
      \end{array}
    \right]\xrightarrow{R_2+(-3)R_1\rightarrow R_2}
    \left[
      \begin{array}{ccc|c}
        1&3&4&7\\
        0&0&-5&-15
      \end{array}
    \right]
  \]
  Now all nonzero rows have been modified, so we need to zero entries above pivots starting from the right, and then scale the pivots to be $1$.
  \[
    \left[
      \begin{array}{ccc|c}
        1&3&4&7\\
        0&0&-5&-15
      \end{array}
    \right]\xrightarrow{R_1+\frac{4}{5}R_2\rightarrow R_1}
    \left[
      \begin{array}{ccc|c}
        1&3&0&-5\\
        0&0&-5&-15
      \end{array}
    \right]\xrightarrow{\frac{-1}{5}R_2\rightarrow R_2}
    \left[
      \begin{array}{ccc|c}
        1&3&0&-5\\
        0&0&1&3
      \end{array}
    \right]
  \]
  Translating back to equations:
  \[
    \begin{cases}
      x_1+3x_2=-5\\
      x_3=3
    \end{cases}
  \]
  So there are infinitely many solutions.

  \Answer
  Follow the algorithm again. The left column is nonzero, so it is a pivot column. The top entry $3$ is not zero, so we use this for our pivot. Then we create zeroes below the pivot.
  \[
    \left[
      \begin{array}{ccc|c}
        3&-2&4&0\\
        9&-6&12&0\\
        6&-4&8&0
      \end{array}
    \right]\xrightarrow{R_2+(-3)R_1\rightarrow R_2}
    \left[
      \begin{array}{ccc|c}
        3&-2&4&0\\
        0&0&0&0\\
        6&-4&8&0
      \end{array}
    \right]\xrightarrow{R_3+(-2)R_1\rightarrow R_3}
    \left[
      \begin{array}{ccc|c}
        3&-2&4&0\\
        0&0&0&0\\
        0&0&0&0
      \end{array}
    \right]
  \]
  Now all nonzero rows have been modified, so we need to zero entries above pivots starting from the right, and then scale the pivots to be $1$.
  \[
    \left[
      \begin{array}{ccc|c}
        3&-2&4&0\\
        0&0&0&0\\
        0&0&0&0
      \end{array}
    \right]\xrightarrow{\frac{1}{3}R_1\rightarrow R_1}
    \left[
      \begin{array}{ccc|c}
        1&-\frac{2}{3}&\frac{4}{3}&0\\
        0&0&0&0\\
        0&0&0&0
      \end{array}
    \right]
  \]
  So this system is consistent with one basic variable ($x_1$) and two free variables ($x_2,x_3$). General solutions to this system have the form:
  \[
    x_2\left[
      \begin{array}{c}
        \frac{2}{3}\\
        1\\
        0
      \end{array}
    \right]+
    x_3\left[
      \begin{array}{c}
        -\frac{4}{3}\\
        0\\
        1
      \end{array}
    \right]
  \]

  \Answer
  No. Since $\vec{v}=2\vec{u}$, then for any scalars $c_u,c_v$, the entries of $c_u\vec{u}+c_v\vec{v}$ are equal:
  \[c_u\vec{u}+c_v\vec{v}=c_u\vec{u}+2c_v\vec{u}=(c_u+2c_v)\vec{u}\]

  \Answer
  We could try the same approach as the previous problem, but when the vectors get larger and more complicated, it can be difficult to analyze them all. Let's use a more systematic approach, and change this to a row reduction problem.

  We are asking whether there are scalars $c_a,c_b,c_c$ such that $c_a\vec{a}+c_b\vec{b}+c_c\vec{c}=\vec{d}$. This is a system of linear equations with the variables $c_a,c_b,c_c$:
  \[
    \left[
      \begin{array}{ccc|c}
        1&0&5&2\\
        -1&1&-6&-1\\
        0&2&8&6
      \end{array}
    \right]
  \]
  Now we follow the row reduction algorithm. The leftmost column is nonzero, and the top entry $1$ is nonzero, so this is our pivot. 
  \[
    \left[
      \begin{array}{ccc|c}
        1&0&5&2\\
        -1&1&-6&-1\\
        0&2&8&6
      \end{array}
    \right]\xrightarrow{R_2+R_1\rightarrow R_2}
    \left[
      \begin{array}{ccc|c}
        1&0&5&2\\
        0&1&-1&1\\
        0&2&8&6
      \end{array}
    \right]
  \]
  Now the second column is nonzero and the leading entry of the second row is $1$, so this is our second pivot.
  \[
    \left[
      \begin{array}{ccc|c}
        1&0&5&2\\
        0&1&-1&1\\
        0&2&8&6
      \end{array}
    \right]\xrightarrow{R_3+(-2)R_2\rightarrow R_3}
    \left[
      \begin{array}{ccc|c}
        1&0&5&2\\
        0&1&-1&1\\
        0&0&10&4
      \end{array}
    \right]
  \]
  Now we create zeroes above pivot positions, and scale the pivot positions to be $1$.
  \[
    \left[
      \begin{array}{ccc|c}
        1&0&5&2\\
        0&1&-1&1\\
        0&0&10&4
      \end{array}
    \right]\xrightarrow{\frac{1}{10}R_3\rightarrow R_3}
    \left[
      \begin{array}{ccc|c}
        1&0&5&2\\
        0&1&-1&1\\
        0&0&1&\frac{2}{5}
      \end{array}
    \right]\xrightarrow{R_2+R_3\rightarrow R_2}
    \left[
      \begin{array}{ccc|c}
        1&0&5&2\\
        0&1&0&\frac{7}{5}\\
        0&0&1&\frac{2}{5}
      \end{array}
    \right]
  \]
  \[
    \xrightarrow{R_1+(-5)R_3\rightarrow R_1}\left[
      \begin{array}{ccc|c}
        1&0&0&0\\
        0&1&0&\frac{7}{5}\\
        0&0&1&\frac{2}{5}
      \end{array}
    \right]
  \]
  So we have our solution $c_a=0, c_b=\frac{7}{5}, c_c=\frac{2}{5}$. Just to be safe, let's check that this actually works:
  \[
    0\left[
      \begin{array}{c}
        1\\
        -1\\
        0
      \end{array}
    \right]+
    \frac{7}{5}\left[
      \begin{array}{c}
        0\\
        1\\
        2
      \end{array}
    \right]+
    \frac{2}{5}\left[
      \begin{array}{c}
        5\\
        -6\\
        8
      \end{array}
    \right]
    =\left[
      \begin{array}{c}
        2\\
        -1\\
        6
      \end{array}
    \right]=\vec{d}
  \]
  \Answer
  \begin{enumerate}
  \item Row reduce the augmented matrix:
    \[
      \left[
        \begin{array}{cc|c}
          1&3&-1\\
          2&4&4
        \end{array}
      \right]\xrightarrow{R_2+(-2)R_1\rightarrow R_2}
      \left[
        \begin{array}{cc|c}
          1&3&-1\\
          0&-2&6
        \end{array}
      \right]\xrightarrow{-\frac{1}{2}R_2\rightarrow R_2}
      \left[
        \begin{array}{cc|c}
          1&3&-1\\
          0&1&-3
        \end{array}
      \right]
    \]
    \[
      \xrightarrow{R_1+(-3)R_2\rightarrow R_1}\left[
        \begin{array}{cc|c}
          1&0&8\\
          0&1&-3
        \end{array}
      \right]      
    \]
    so $\vec{b}\in \Span\{\vec{u},\vec{v}\}$ (Check: $8\left[\begin{array}{c} 1\\2\end{array}\right]+(-3)\left[\begin{array}{c} 3\\4\end{array}\right]=\left[\begin{array}{c} -1\\4\end{array}\right]=\vec{b}$).
    
  \item If $\vec{b}\notin \mathbb{R}^2$, then $\vec{b}$ is not in the span of $\vec{u}$ and $\vec{v}$. This is a somewhat silly answer to this question. Now assume $\vec{b}=\left[\begin{array}{c} x\\y\end{array}\right] \in \mathbb{R}^2$. Then we row reduce the augmented matrix:
    \[
      \left[
        \begin{array}{cc|c}
          1&3&x\\
          2&4&y
        \end{array}
      \right]\xrightarrow{R_2+(-2)R_1\rightarrow R_2}
      \left[
        \begin{array}{cc|c}
          1&3&x\\
          0&-2&y-2x
        \end{array}
      \right]\xrightarrow{-\frac{1}{2}R_2\rightarrow R_2}
      \left[
        \begin{array}{cc|c}
          1&3&x\\
          0&1&\frac{2x-y}{2}
        \end{array}
      \right]
    \]
    \[
      \xrightarrow{R_1+(-3)R_2\rightarrow R_1}\left[
        \begin{array}{cc|c}
          1&0&-2x+\frac{3}{2}y\\
          0&1&\frac{2x-y}{2}
        \end{array}
      \right]      
    \]
    so $\vec{b}\in\Span\{\vec{u},\vec{v}\}$. (Check $(-2x+\frac{3}{2}y)\left[\begin{array}{c} 1\\2\end{array}\right]+\frac{2x-y}{2}\left[\begin{array}{c} 3\\4\end{array}\right]=\left[\begin{array}{c} x\\y\end{array}\right]=\vec{b}$.)
    
  \end{enumerate}
  
\end{enumerate}
\end{comment}
%\end{document}


%%% Local Variables: 
%%% mode: latex
%%% TeX-master: t
%%% End: 
